\chapter{Centroid Decomposition}

Heavy-light decomposition, in the previous chapter, took a tree and sliced it into
chains so that root-to-vertex paths cross only $\Oh(\log n)$ of them. That is the
right hammer when your queries follow the tree's own paths -- update a vertex, ask
for the maximum on the path from $u$ to $v$. But a whole family of tree problems is
not about one path at a time; it is about \emph{all} paths at once. How many pairs
of vertices are at distance exactly $k$? For each vertex, what is the closest red
vertex anywhere in the tree? Add a vertex to a set; report the nearest set-member
to a query vertex. These ask you to reason about every path simultaneously, and a
chain decomposition has nothing to say about them.

Centroid decomposition is the hammer for those. The idea is simple:
find a vertex whose removal splits the tree into pieces each no larger than half
the tree -- a \emph{centroid} -- delete it, and recurse on the pieces. Because
every split at least halves the part you recurse into, the recursion is only
$\Oh(\log n)$ deep. The centroids, linked parent-to-child by the recursion, form a
new tree of logarithmic height: the \emph{centroid tree}. Its decisive property is
that \emph{every} path in the original tree passes through the centroid-tree
ancestor of its two endpoints. That single fact converts ``reason about all
$\binom{n}{2}$ paths'' into ``for each of $n$ centroids, reason about the paths
that pass through it,'' and the halving guarantees the bookkeeping stays cheap.

We will define centroids and prove they always exist, build the decomposition and
prove its depth bound, then spend the heart of the chapter on the technique that
makes it worth knowing: solving path and distance problems by attributing each path
to the unique centroid that owns it. A small tree is traced by hand throughout.

\chapterline

\section{Tree centroids}
\label{ch20:centroids}

Throughout, $T=(V,E)$ is a tree on $n=\abs{V}$ vertices: connected, acyclic,
$n-1$ edges. Deleting a vertex $v$ from a tree disconnects it into one component
per neighbor of $v$; we call these the \emph{branches} hanging off $v$, and write
their sizes as $s_1,\dots,s_{\deg v}$ with $\sum_i s_i = n-1$.

\begin{definition}[Centroid]
\label{ch20:def-centroid}
A vertex $c\in V$ is a \emph{centroid} of $T$ if removing $c$ leaves every
resulting component with at most $\floor{n/2}$ vertices. Equivalently, every
branch hanging off $c$ has size at most $\floor{n/2}$.
\end{definition}

\begin{intuitionbox}
A centroid is the tree's balance point. Stand on a vertex and look down each
branch; if one branch holds more than half of everything, you are off-center --
step toward that heavy branch and you improve. A centroid is the place where no
single step toward any branch helps, because no branch is heavy enough to be worth
walking into. Contrast the \emph{root} of a tree, which you pick arbitrarily and
which can be wildly unbalanced (a path rooted at its end). The centroid is forced
by the tree's shape, not chosen.
\end{intuitionbox}

Here is the small tree we will use for the entire chapter, shown in Figure~\ref{ch20:fig-running-tree}. It has $n=9$ vertices.

\begin{figure}[H]
\centering
\begin{tikzpicture}[every node/.style={circle,draw=nodeblue,thick,fill=blue!5,minimum size=7mm,
  font=\small},level distance=12mm,sibling distance=16mm]
\node (n4) at (0,0) {4}
  child {node (n2) {2}
    child {node (n1) {1}}
    child {node (n3) {3}}
  }
  child {node[draw=structurepurple,very thick,fill=purple!12] (n6) {6}
    child {node (n5) {5}}
    child {node (n7) {7}
      child {node (n8) {8}}
      child {node (n9) {9}}
    }
  };
\node[draw=none,fill=none,structurepurple,font=\footnotesize,right=3mm of n6] {centroid $c=6$};
\end{tikzpicture}
\caption{The running 9-vertex tree. Removing node $4$ leaves a heavy component of size $5$ ($\{5,6,7,8,9\}$). Stepping toward the heavy branch into node $6$ yields three components of sizes $4$ ($\{1,2,3,4\}$), $1$ ($\{5\}$), and $3$ ($\{7,8,9\}$)---each at most $\floor{9/2}=4$. Thus vertex $6$ is the centroid.}
\label{ch20:fig-running-tree}
\end{figure}

The edge set, pinned down so the hand-traces are unambiguous: $4$ is adjacent to
$2$ and $6$; $2$ is adjacent to $1,3,4$; $6$ is adjacent to $4,5,7$; $7$ is
adjacent to $6,8,9$.

Let us locate a centroid by hand. Total $n=9$, so a branch is ``heavy'' if it
holds more than $\floor{9/2}=4$ vertices. Try vertex $4$: removing it leaves the
component $\{2,1,3\}$ of size $3$ and the component $\{6,5,7,8,9\}$ of size $5$.
The second branch has $5>4$ vertices, so $4$ is \emph{not} a centroid; we should
step toward the heavy side, into $6$. Remove $6$: components are $\{4,2,1,3\}$
(size $4$), $\{5\}$ (size $1$), $\{7,8,9\}$ (size $3$). Every piece is $\le 4$, so
\textbf{vertex $6$ is the centroid}.

\begin{whybox}
The ``walk toward the heavy branch'' instinct is exactly the existence proof.
Define a potential $\Phi(v)=$ size of the largest branch hanging off $v$. If $v$
is not a centroid, some branch $B$ has $\abs{B}>\floor{n/2}\ge (n-1)/2$, hence
$\abs{B}\ge \floor{n/2}+1 > n - \abs{B}$, so $B$ holds a strict majority. Step
from $v$ to its neighbor $w$ inside $B$. From $w$'s vantage, the branch back
toward $v$ contains everything outside $B$ plus $v$ itself: $n-\abs{B} < \abs{B}$
vertices, strictly fewer than $\Phi(v)$. The other branches off $w$ are sub-pieces
of $B\setminus\{w\}$, each strictly smaller than $\abs{B}=\Phi(v)$. So
$\Phi(w)<\Phi(v)$. The potential is a positive integer that strictly decreases
each step, so the walk cannot loop and must halt -- and it can only halt at a
centroid.
\end{whybox}

\begin{theorem}[Existence of a centroid]
\label{ch20:thm-exists}
Every tree on $n\ge 1$ vertices has at least one centroid. Moreover it has at most
two, and if it has two they are adjacent.
\end{theorem}

\begin{proof}
\emph{Existence.} The walk in the box above starts at any vertex and strictly
decreases the integer potential $\Phi$ each step, so after finitely many steps it
reaches a vertex with no heavy branch -- a centroid.

\emph{At most two, and adjacent.} Suppose $c$ is a centroid and $c'\neq c$ is also
a centroid. Root the tree at $c$ and let $B$ be the branch off $c$ that contains
$c'$. Because $c$ is a centroid, $\abs{B}\le\floor{n/2}$. Now view things from
$c'$: one of $c'$'s branches is the part of the tree on $c$'s side, which has size
$n-\abs{\mathrm{subtree}(c')}$ where the subtree is taken in the $c$-rooted tree.
Since $c'$ is a centroid this side has size $\le\floor{n/2}$, giving
$\abs{\mathrm{subtree}(c')}\ge n-\floor{n/2}=\ceil{n/2}$. But
$\mathrm{subtree}(c')\subseteq B$, so $\ceil{n/2}\le\abs{B}\le\floor{n/2}$, forcing
$\abs{B}=\ceil{n/2}=\floor{n/2}$, which needs $n$ even and pins the subtree of
$c'$ to be exactly $B$. Then the $c$-to-$c'$ branch has size exactly $n/2$ on both
ends, and any vertex strictly between $c$ and $c'$ on the path would push the
larger side past $n/2$ -- contradiction unless $c$ and $c'$ are adjacent.
\end{proof}

\begin{trapbox}
The centroid is not the same as the \emph{center} (the vertex minimizing
eccentricity, i.e.\ the middle of a longest path), nor the \emph{barycenter} of a
weighted embedding. On the path graph $1\!-\!2\!-\!3\!-\!4\!-\!5$ all three
coincide at $3$, which lulls people into thinking they are interchangeable. They
are not: take a star with one long whisker and the center and centroid drift apart.
For centroid decomposition you specifically need the size-balance property of
Definition~\ref{ch20:def-centroid}; eccentricity is irrelevant.
\end{trapbox}

\subsection{Finding a centroid in linear time}

The walk-toward-the-heavy-branch argument is constructive but a little awkward to
code directly. The clean implementation roots the (sub)tree arbitrarily, computes
subtree sizes with one DFS, then in a second pass walks down from the root always
stepping into a child whose subtree exceeds half, stopping when none does.

\begin{algorithm}[H]
\caption{Find the centroid of a tree (or subtree) of size $m$}
\label{ch20:alg-find}
\DontPrintSemicolon
\KwIn{Tree given by adjacency lists; a start vertex $r$; a $\mathrm{removed}[\cdot]$
flag marking deleted vertices; the count $m$ of live vertices reachable from $r$.}
\KwOut{A centroid of that component.}
\SetKwProg{Fn}{procedure}{}{}
\Fn{\textsc{SubtreeSizes}$(u,p)$}{
  $\mathrm{sz}[u]\gets 1$\;
  \ForEach{$w\in\mathrm{adj}(u)$ with $w\neq p$ and not $\mathrm{removed}[w]$}{
    $\mathrm{sz}[u]\gets\mathrm{sz}[u]+\textsc{SubtreeSizes}(w,u)$\;
  }
  \Return $\mathrm{sz}[u]$\;
}
\Fn{\textsc{FindCentroid}$(u,p,m)$}{
  \ForEach{$w\in\mathrm{adj}(u)$ with $w\neq p$ and not $\mathrm{removed}[w]$}{
    \If{$\mathrm{sz}[w] > m/2$}{
      \Return \textsc{FindCentroid}$(w,u,m)$\;
    }
  }
  \Return $u$\;
}
$\textsc{SubtreeSizes}(r,\bot)$\;
\Return $\textsc{FindCentroid}(r,\bot,m)$\;
\end{algorithm}

\begin{whybox}
After \textsc{SubtreeSizes}$(r,\cdot)$, $\mathrm{sz}[w]$ is the size of $w$'s
subtree in the $r$-rooted component. At a vertex $u$, the branch hanging off $u$
toward parent $p$ has size $m-\mathrm{sz}[u]$, and each child-branch has size
$\mathrm{sz}[w]$. \textsc{FindCentroid} descends only into a child whose subtree
$>m/2$. Two facts make this correct. First, at most one child can have
$\mathrm{sz}[w]>m/2$, since two such would already overflow $m$; so the descent is
deterministic. Second, when we step from $u$ into such a $w$, the new
parent-branch (back toward $u$) has size $m-\mathrm{sz}[w]<m/2$, so it can never
again be the heavy branch -- we only ever move ``inward'' and sizes shrink.
Termination at $u$ means no child exceeds $m/2$ and the parent-branch
$m-\mathrm{sz}[u]\le m/2$, so every branch off $u$ is $\le\floor{m/2}$ and $u$ is a
centroid. Both passes are a single DFS over the $m$ live vertices: $\Oh(m)$ time.
\end{whybox}

\section{The decomposition and its depth}
\label{ch20:decomp}

Now recurse. Find a centroid $c$ of the whole tree, record it as the root of the
centroid tree, mark it removed, and recurse on each remaining component; the
centroid found in a component becomes a \emph{child} of $c$ in the centroid tree.

\begin{algorithm}[H]
\caption{Build the centroid decomposition}
\label{ch20:alg-build}
\DontPrintSemicolon
\KwIn{Tree $T=(V,E)$ as adjacency lists.}
\KwOut{$\mathrm{par}[\cdot]$: parent of each vertex in the centroid tree.}
\SetKwProg{Fn}{procedure}{}{}
\ForEach{$v\in V$}{$\mathrm{removed}[v]\gets\textsc{false}$}
\Fn{\textsc{Decompose}$(r,\mathrm{cpar})$}{
  $m\gets$ number of live vertices in $r$'s component (one DFS)\;
  $c\gets\textsc{FindCentroid}(r,\bot,m)$ \tcp*{Algorithm~\ref{ch20:alg-find}}
  $\mathrm{par}[c]\gets\mathrm{cpar}$\;
  $\mathrm{removed}[c]\gets\textsc{true}$\;
  \ForEach{$w\in\mathrm{adj}(c)$ with not $\mathrm{removed}[w]$}{
    \textsc{Decompose}$(w,c)$\;
  }
}
\textsc{Decompose}$(\text{any vertex},\ \bot)$\;
\Return $\mathrm{par}$\;
\end{algorithm}

\subsection{Tracing it on the nine-vertex tree}

\begin{labbox}
Run Algorithm~\ref{ch20:alg-build} on our tree. We found the top centroid is $6$.
Remove $6$; three live components remain.
\begin{itemize}
  \item Component $\{4,2,1,3\}$, $m=4$, ``heavy'' threshold $>2$. Root at $4$:
        subtree of $2$ is $\{2,1,3\}$, size $3>2$, step into $2$; from $2$ the
        parent-branch is $\{4\}$ size $1$, children $\{1\},\{3\}$ size $1$ each --
        none $>2$. \textbf{Centroid $=2$}, child of $6$. Recurse: removing $2$
        gives $\{4\},\{1\},\{3\}$, each a single vertex and its own centroid,
        children of $2$.
  \item Component $\{5\}$, $m=1$. \textbf{Centroid $=5$}, child of $6$ (a leaf).
  \item Component $\{7,8,9\}$, $m=3$, threshold $>1$. Root at $7$: children
        $\{8\},\{9\}$ size $1$, parent-branch empty. \textbf{Centroid $=7$}, child
        of $6$. Recurse: $\{8\},\{9\}$ become children of $7$.
\end{itemize}
\end{labbox}

The recursive component splitting process and the resulting balanced centroid tree are shown in Figure~\ref{ch20:fig-centroid-hierarchy}.

\begin{figure}[H]
\centering
\begin{tikzpicture}[
  every node/.style={font=\small},
  cnode/.style={circle,draw=structurepurple,very thick,fill=purple!15,minimum size=6.5mm,inner sep=0pt},
  vnode/.style={circle,draw=nodeblue,thick,fill=blue!5,minimum size=6mm,inner sep=0pt},
  compbox/.style={draw=black!25,rounded corners=4pt,fill=black!2,inner sep=3pt},
  splitarrow/.style={-{Stealth[length=2mm]},edgegray,thick}
]
% Panel (a): Recursive Component Splitting
\begin{scope}
  \node[font=\footnotesize\bfseries,align=center] at (3.2, 3.8) {\textbf{(a)} Recursive Tree Partitioning};

  % Level 0
  \node[draw=black!30,rounded corners=4pt,fill=purple!4,inner sep=4pt] (lvl0) at (3.2, 2.8) {
    \begin{tikzpicture}[baseline]
      \node[vnode] (a4) at (0,0.6) {4};
      \node[vnode] (a2) at (-0.6,0) {2};
      \node[cnode] (a6) at (0.6,0) {6};
      \draw[thick] (a4)--(a2); \draw[thick] (a4)--(a6);
      \node[font=\scriptsize\color{structurepurple},right=1pt of a6] {$\leftarrow$ root centroid $c_0=6$};
    \end{tikzpicture}
  };

  % Down arrows to Level 1
  \draw[splitarrow] (2.2, 2.1) -- (1.0, 1.4);
  \draw[splitarrow] (3.2, 2.1) -- (3.2, 1.4);
  \draw[splitarrow] (4.2, 2.1) -- (5.4, 1.4);

  % Level 1 Components
  \node[compbox] (c1) at (1.0, 0.8) {
    \begin{tikzpicture}[baseline]
      \node[cnode] (b2) at (0,0.4) {2};
      \node[vnode] (b4) at (-0.4,-0.2) {4};
      \node[vnode] (b1) at (0,-0.2) {1};
      \node[vnode] (b3) at (0.4,-0.2) {3};
      \draw (b2)--(b4); \draw (b2)--(b1); \draw (b2)--(b3);
    \end{tikzpicture}
  };
  \node[font=\tiny,below=1pt of c1] {$\{1,2,3,4\}$, $c_1=2$};

  \node[compbox] (c2) at (3.2, 0.8) {
    \begin{tikzpicture}[baseline]
      \node[cnode] (b5) at (0,0) {5};
    \end{tikzpicture}
  };
  \node[font=\tiny,below=1pt of c2] {$\{5\}$, $c=5$};

  \node[compbox] (c3) at (5.4, 0.8) {
    \begin{tikzpicture}[baseline]
      \node[cnode] (b7) at (0,0.3) {7};
      \node[vnode] (b8) at (-0.3,-0.2) {8};
      \node[vnode] (b9) at (0.3,-0.2) {9};
      \draw (b7)--(b8); \draw (b7)--(b9);
    \end{tikzpicture}
  };
  \node[font=\tiny,below=1pt of c3] {$\{7,8,9\}$, $c_2=7$};

  % Down arrows to Level 2
  \draw[splitarrow] (1.0, 0.2) -- (1.0, -0.5);
  \draw[splitarrow] (5.4, 0.2) -- (5.4, -0.5);

  % Level 2 Singletons
  \node[font=\scriptsize,edgegray] at (1.0, -0.8) {Singletons: $\{1\}, \{3\}, \{4\}$};
  \node[font=\scriptsize,edgegray] at (5.4, -0.8) {Singletons: $\{8\}, \{9\}$};

  \node[font=\footnotesize,align=center] at (3.2, -1.6)
    {Each step removes a centroid and cuts component\\
     sizes to $\le m/2$; depth $\le \floor{\log_2 9} = 3$};
\end{scope}

% Panel (b): Resulting Centroid Tree
\begin{scope}[xshift=7.8cm]
  \node[font=\footnotesize\bfseries,align=center] at (2.2, 3.8) {\textbf{(b)} Balanced Centroid Tree $\mathcal{T}_C$};

  \node[cnode] (r6) at (2.2, 2.8) {6};

  \node[cnode] (r2) at (0.7, 1.4) {2};
  \node[cnode] (r5) at (2.2, 1.4) {5};
  \node[cnode] (r7) at (3.7, 1.4) {7};

  \node[vnode] (r4) at (0.0, 0.0) {4};
  \node[vnode] (r1) at (0.7, 0.0) {1};
  \node[vnode] (r3) at (1.4, 0.0) {3};

  \node[vnode] (r8) at (3.2, 0.0) {8};
  \node[vnode] (r9) at (4.2, 0.0) {9};

  \draw[thick,structurepurple] (r6) -- (r2);
  \draw[thick,structurepurple] (r6) -- (r5);
  \draw[thick,structurepurple] (r6) -- (r7);

  \draw[thick,structurepurple] (r2) -- (r4);
  \draw[thick,structurepurple] (r2) -- (r1);
  \draw[thick,structurepurple] (r2) -- (r3);

  \draw[thick,structurepurple] (r7) -- (r8);
  \draw[thick,structurepurple] (r7) -- (r9);

  \node[font=\scriptsize,edgegray,right=2pt of r6] {depth $0$ ($m \le 9$)};
  \node[font=\scriptsize,edgegray,right=2pt of r7] {depth $1$ ($m \le 4$)};
  \node[font=\scriptsize,edgegray,right=2pt of r9] {depth $2$ ($m \le 2$)};

  \node[font=\footnotesize,align=center] at (2.2, -1.6)
    {Height $= 2 \le \floor{\log_2 9}$;\\
     every path $u \rightsquigarrow v$ passes through\\
     $\mathrm{LCA}_{\mathcal{T}_C}(u, v)$};
\end{scope}
\end{tikzpicture}
\caption{The centroid tree recursion hierarchy. (a)~Recursive decomposition:
at each stage, removing the centroid $c$ guarantees that all disconnected components
have at most $\floor{m/2}$ vertices, driving logarithmic recursion depth.
(b)~The resulting centroid tree $\mathcal{T}_C$ has height $2 \le \log_2 n$.
The unique lowest common ancestor $a = \mathrm{LCA}_{\mathcal{T}_C}(u, v)$ in the
centroid tree is the unique centroid that separates and owns the path between $u$ and $v$.}
\label{ch20:fig-centroid-hierarchy}
\end{figure}

Notice this is a different tree on the same nine vertices. It is broad and shallow:
height $2$ (edges $6\to2\to4$). The original tree had height $3$. In general the
centroid tree's height is $\Oh(\log n)$, which we now prove.

\begin{theorem}[Logarithmic depth]
\label{ch20:thm-depth}
The centroid tree produced by Algorithm~\ref{ch20:alg-build} has height at most
$\floor{\log_2 n}$. Consequently every vertex has at most $\floor{\log_2 n}+1$
ancestors in the centroid tree (itself included).
\end{theorem}

\begin{proof}
When \textsc{Decompose} processes a component of size $m$ and removes its centroid
$c$, every child component has size $\le\floor{m/2}$ by
Definition~\ref{ch20:def-centroid}. The component handled at centroid-tree depth
$d$ (root has depth $0$) therefore has size at most $\floor{n/2^{\,d}}\le n/2^{\,d}$.
A component must contain at least one vertex to exist, so $n/2^{\,d}\ge 1$, i.e.\
$d\le\log_2 n$. Depth is an integer, so $d\le\floor{\log_2 n}$.
\end{proof}

\begin{complexitybox}
\textbf{Build time.} Let $T(n)$ be the work to decompose a component of size $n$.
Finding the centroid and counting the component is $\Oh(n)$ (one DFS), and the
children's sizes sum to $\le n-1$. So the total work at each \emph{level} of the
centroid tree is $\Oh(\text{sum of component sizes at that level})=\Oh(n)$, and
there are $\Oh(\log n)$ levels (Theorem~\ref{ch20:thm-depth}). Hence
\[
  T(n)=\Oh(n\log n)\quad\text{time},\qquad \Oh(n)\ \text{extra space}.
\]
The recurrence $T(n)\le \Oh(n) + \sum_i T(s_i)$ with $\sum_i s_i\le n-1$ and each
$s_i\le n/2$ solves to the same bound. The centroid tree itself stores one parent
pointer per vertex: $\Oh(n)$.
\end{complexitybox}

\begin{trapbox}
A seductive bug: computing component sizes by DFS-ing the \emph{whole original
tree} each time instead of only the live component. If \textsc{SubtreeSizes}
ignores $\mathrm{removed}$ and walks the full tree, every call is $\Oh(n)$ and you
do $n$ of them in the worst case (a path), so build degrades to $\Oh(n^2)$. The
$\Oh(n\log n)$ bound relies on each DFS touching only \emph{currently live}
vertices, which is why the $\mathrm{removed}[\cdot]$ guard appears in every loop of
Algorithm~\ref{ch20:alg-find}. Always re-measure the local component; never the
global tree.
\end{trapbox}

\begin{historybox}
Centroids of trees go back to Camille Jordan, who in 1869 proved every tree has a
center and a centroid while studying the symmetry of molecular trees in chemistry
-- the same Jordan of curve-theorem and normal-form fame. The recursive
\emph{decomposition} for algorithmic divide-and-conquer is folklore from the
1980s--1990s separator literature (centroids are the simplest tree separators),
and it became a standard competitive-programming technique in the 2000s for
path-counting and nearest-marked-vertex queries. It is the tree special case of
the general graph-separator paradigm you will meet again in the planar-separator
and treewidth chapters.
\end{historybox}

\section{Solving path problems by decomposition}
\label{ch20:pathkey}

Everything above is scaffolding for one structural lemma, which is the reason the
technique exists.

\begin{lemma}[Path--centroid lemma]
\label{ch20:lem-pathkey}
Let $u,v\in V$ and let $a$ be their lowest common ancestor in the centroid tree.
Then the unique tree path $u\!\rightsquigarrow\! v$ in $T$ passes through $a$, and
$a$ is the \emph{only} vertex on that path that is an ancestor of both $u$ and $v$
in the centroid tree.
\end{lemma}

\begin{proof}
Consider the moment \textsc{Decompose} acts on the component that contains both $u$
and $v$ and is about to remove its centroid. While $a=\mathrm{lca}(u,v)$ has not
yet been removed, $u$ and $v$ live in the same component (that is what being
descendants of the same centroid-tree node means). At the step that removes $a$,
either $a\in\{u,v\}$, or $a$ splits this component so that $u$ and $v$ fall into
\emph{different} child components -- otherwise their lowest common ancestor would
lie deeper, inside that single child, contradicting $a=\mathrm{lca}(u,v)$. In a
tree, if deleting $a$ separates $u$ from $v$, then $a$ lies on the path
$u\!\rightsquigarrow\! v$ (the path is the only $u$--$v$ route, and it is cut). So
$a$ is on the path.

Uniqueness: suppose some other vertex $b\neq a$ on the path were a common ancestor
of $u$ and $v$ in the centroid tree. A common ancestor of $u$ and $v$ is, in
particular, an ancestor of $a=\mathrm{lca}(u,v)$ or equals $a$; if $b\neq a$ then
$b$ is a strict ancestor of $a$, so $b$ was removed at a shallower centroid level,
\emph{before} $a$. But at the moment $b$ was removed, $u$ and $v$ still lay in the
same component (they are removed even later), so $b$ did not separate them, hence
$b$ is not on the path $u\!\rightsquigarrow\! v$ -- contradiction. Thus $a$ is the
only such vertex.
\end{proof}

\begin{intuitionbox}
Read the lemma as an \emph{ownership} rule. Every one of the $\binom{n}{2}$ paths
in $T$ is owned by exactly one centroid: the highest centroid that the path
touches, equivalently $\mathrm{lca}$ of its endpoints in the centroid tree. If you
want to process all paths, process each centroid $c$ once and handle only the paths
that pass through $c$ and live entirely within $c$'s component. No path is missed
(every path has an owner) and no path is double-counted (the owner is unique). The
$\Oh(\log n)$ depth means each vertex belongs to only $\Oh(\log n)$ components, so
the total ``path through $c$'' work, summed over all centroids, stays small.
\end{intuitionbox}

The standard recipe. For each centroid $c$, every relevant path through $c$ is the
concatenation of two ``arms'': a path from $c$ down into one branch, and a path
from $c$ down into another (or the trivial empty arm ending at $c$). Distances
along an arm are just distances to $c$ inside $c$'s component, which one BFS/DFS
from $c$ computes. The problem then becomes: combine arms from \emph{different}
branches. Two classic combination patterns follow.

\begin{figure}[H]
\centering
\begin{tikzpicture}[every node/.style={circle,draw,nodeblue,thick,
  minimum size=7mm,font=\small}]
\node (c) at (0,0) {$c$};
\node (u) at (-2.2,-1.6) {$u$};
\node (v) at (2.2,-1.6) {$v$};
\draw[edgegray,thick,dashed] (c) -- node[draw=none,fill=none,above left,font=\scriptsize]{$d_u$} (u);
\draw[edgegray,thick,dashed] (c) -- node[draw=none,fill=none,above right,font=\scriptsize]{$d_v$} (v);
\begin{scope}[on background layer]
  \fill[blue!6] (-2.2,-1.6) ++(-0.9,-0.5) -- ++(1.8,0) -- ++(-0.9,2.0) -- cycle;
  \fill[green!6] (2.2,-1.6) ++(-0.9,-0.5) -- ++(1.8,0) -- ++(-0.9,2.0) -- cycle;
\end{scope}
\node[draw=none,fill=none,font=\scriptsize,blue!50!black] at (-2.2,-2.4) {branch $B_1$};
\node[draw=none,fill=none,font=\scriptsize,green!40!black] at (2.2,-2.4) {branch $B_2$};
\end{tikzpicture}
\caption{A path owned by centroid $c$ splits into two arms in distinct branches:
$d_u=\dist(c,u)$ and $d_v=\dist(c,v)$, with total length $d_u+d_v$. Two arms in
the \emph{same} branch do not form a path through $c$ and must be excluded by the
per-branch correction.}
\label{ch20:fig-arms}
\end{figure}

\subsection{Pattern A: counting paths of a given length}

Problem: count unordered pairs $\{u,v\}$ with $\dist_T(u,v)=k$ (here $\dist$ is the
number of edges). Each such pair is owned by one centroid $a=\mathrm{lca}(u,v)$,
and the path splits at $a$ into two arms of lengths $d_u=\dist(a,u)$ and
$d_v=\dist(a,v)$ with $d_u+d_v=k$, the two arms living in different branches off
$a$ (or one arm empty when $a\in\{u,v\}$); see Figure~\ref{ch20:fig-arms}.

So at centroid $c$: run a DFS to get, for every live vertex $x$, its depth
$\dist(c,x)$ and which branch of $c$ it lies in. Build a histogram
$\mathrm{cnt}[d]=$ number of live vertices at depth $d$ from $c$. The number of
\emph{ordered} arm pairs summing to $k$ is $\sum_{d} \mathrm{cnt}[d]\cdot
\mathrm{cnt}[k-d]$, but this overcounts: it includes pairs of arms in the
\emph{same} branch (whose true path does not pass through $c$). Subtract them with
the standard inclusion--exclusion: for each branch $B$, build $\mathrm{cnt}_B[\cdot]$
and subtract $\sum_d \mathrm{cnt}_B[d]\cdot\mathrm{cnt}_B[k-d]$. Finally divide the
ordered count by two.

\begin{algorithm}[H]
\caption{Count pairs at distance exactly $k$ via centroid decomposition}
\label{ch20:alg-count}
\DontPrintSemicolon
\KwIn{Tree $T$; target distance $k$.}
\KwOut{$\mathrm{ans}=\abs{\{\{u,v\}:\dist_T(u,v)=k\}}$.}
$\mathrm{ans}\gets 0$\;
\SetKwProg{Fn}{procedure}{}{}
\Fn{\textsc{CountPaths}$(\mathrm{cnt}[\cdot],\mathrm{sign})$}{
  \tcp{add $\mathrm{sign}\cdot(\text{ordered arm pairs summing to }k)$}
  \For{$d\gets 0$ \KwTo $k$}{
    $\mathrm{ans}\gets\mathrm{ans}+\mathrm{sign}\cdot\mathrm{cnt}[d]\cdot\mathrm{cnt}[k-d]$\;
  }
}
\Fn{\textsc{Solve}$(r)$}{
  $m\gets$ size of $r$'s live component\;
  $c\gets\textsc{FindCentroid}(r,\bot,m)$; $\mathrm{removed}[c]\gets\textsc{true}$\;
  $\mathrm{cnt}\gets$ depth histogram of whole component from $c$ ($\mathrm{cnt}[0]=1$ for $c$)\;
  \textsc{CountPaths}$(\mathrm{cnt},+1)$\;
  \ForEach{branch $B$ off $c$ (each live neighbor's subtree)}{
    $\mathrm{cnt}_B\gets$ depth histogram of $B$ from $c$ (depths $\ge 1$)\;
    \textsc{CountPaths}$(\mathrm{cnt}_B,-1)$\;
  }
  \ForEach{$w\in\mathrm{adj}(c)$ with not $\mathrm{removed}[w]$}{\textsc{Solve}$(w)$\;}
}
\textsc{Solve}$(\text{any vertex})$\;
\Return $\mathrm{ans}/2$\;
\end{algorithm}

\begin{whybox}
The whole-component histogram counts every ordered pair of arms summing to $k$,
including illegal same-branch pairs and including the $\mathrm{cnt}[0]=1$ slot for
$c$ itself. Subtracting each branch's self-pairs removes exactly the same-branch
pairs: two arms in branch $B$ both contribute to $\mathrm{cnt}_B$ and get
cancelled. What survives is cross-branch ordered pairs, plus the two ordered pairs
$(c,x)$ and $(x,c)$ for each $x$ at depth $k$ -- those use the $d=0$ slot, which is
present only in the full histogram (each $\mathrm{cnt}_B$ starts at depth $1$). So a
distance-$k$ path with one endpoint $=c$ is also counted twice. Each legitimate
path through $c$ is thus counted exactly twice, and the final division by $2$ is
exact. By Lemma~\ref{ch20:lem-pathkey} every distance-$k$ pair is owned by exactly
one centroid, so summing over centroids counts each pair exactly once.
\end{whybox}

\begin{complexitybox}
At centroid $c$ the work is $\Oh(m_c)$ for the DFS plus $\Oh(\min(k,m_c))$ per
histogram convolution; summed naively the convolutions cost $\Oh(\sum_c k)$, which
can be $\Oh(nk)$. Capping each histogram at the branch's max depth makes total
convolution work $\Oh(\sum_c m_c)=\Oh(n\log n)$, since each vertex appears in
$\Oh(\log n)$ components and contributes $\Oh(1)$ to one histogram slot per
component. Overall $\Oh(n\log n)$ time, $\Oh(n)$ space -- versus the $\Oh(n^2)$ of
checking all pairs.
\end{complexitybox}

\begin{labbox}
Trace Algorithm~\ref{ch20:alg-count} with $k=2$ at the top centroid $c=6$. DFS from
$6$ over the whole tree gives depths $\dist(6,\cdot)$: $6\!:\!0$; $4,5,7\!:\!1$;
$2,8,9\!:\!2$; $1,3\!:\!3$. Histogram $\mathrm{cnt}=[1,3,3,2]$ for depths
$0,1,2,3$. Ordered pairs summing to $k=2$ from the full histogram:
$\mathrm{cnt}[0]\mathrm{cnt}[2]+\mathrm{cnt}[1]\mathrm{cnt}[1]+\mathrm{cnt}[2]\mathrm{cnt}[0]
=3+9+3=15$. Now the three branches off $6$. Branch $\{4,2,1,3\}$ has depths
$4\!:\!1,\,2\!:\!2,\,1\!:\!3,\,3\!:\!3$, so $\mathrm{cnt}_B$ has $1$ at depth $1$,
$1$ at depth $2$, $2$ at depth $3$; ordered pairs summing to $2$:
$\mathrm{cnt}_B[1]^2=1$. Branch $\{5\}$: depth $1$, gives $1$. Branch $\{7,8,9\}$:
depths $7\!:\!1,\,8\!:\!2,\,9\!:\!2$, gives $\mathrm{cnt}_B[1]^2=1$. Subtract
$1+1+1=3$, leaving $15-3=12$ ordered pairs owned by $6$, i.e.\ $6$ unordered
distance-$2$ pairs through $6$: the three ``via $6$'' pairs $\{4,5\},\{4,7\},\{5,7\}$
and the three depth-$2$ endpoints-with-$6$ pairs $\{6,2\},\{6,8\},\{6,9\}$.
The remaining distance-$2$ pairs ($\{1,3\}$ via $2$, $\{8,9\}$ via $7$) are owned by
the deeper centroids $2$ and $7$ and are counted when \textsc{Solve} recurses
there. Summing each centroid's contribution and dividing by $2$ gives the total.
\end{labbox}

\subsection{Pattern B: nearest marked vertex under updates}

Problem: maintain a set $S\subseteq V$ (initially empty); support
\textsc{Mark}$(v)$ adding $v$ to $S$, and \textsc{Query}$(v)$ returning
$\min_{s\in S}\dist_T(v,s)$. This is the canonical \emph{online} centroid-tree
application and shows off the data structure rather than a one-shot count.

Store at each centroid $c$ a value $\mathrm{best}[c]=\min_{s\in S\cap
\mathrm{comp}(c)}\dist_T(c,s)$, the distance from $c$ to the nearest marked vertex
in $c$'s component (where $\mathrm{comp}(c)$ is the set of vertices for which $c$ is
a centroid-tree ancestor). To support distance lookups we precompute, for every
vertex $x$ and every centroid-tree ancestor $a$ of $x$, the value
$\dist_T(x,a)$ -- there are only $\Oh(\log n)$ ancestors per vertex, so this is an
$n\times\Oh(\log n)$ table built during the decomposition, where each centroid's
BFS over its live component records $\dist_T(x,c)$ for every $x$ it reaches.

\begin{algorithm}[H]
\caption{Nearest marked vertex (Pattern B)}
\label{ch20:alg-nearest}
\DontPrintSemicolon
\SetKwProg{Fn}{procedure}{}{}
\KwData{$\mathrm{par}[\cdot]$ centroid tree; $D[x][a]=\dist_T(x,a)$ for each
centroid-ancestor $a$ of $x$; $\mathrm{best}[c]\gets+\infty$ for all $c$.}
\Fn{\textsc{Mark}$(v)$}{
  $a\gets v$\;
  \While{$a\neq\bot$}{
    $\mathrm{best}[a]\gets\min(\mathrm{best}[a],\,D[v][a])$\;
    $a\gets\mathrm{par}[a]$\;
  }
}
\Fn{\textsc{Query}$(v)$}{
  $a\gets v$; $\mathrm{res}\gets+\infty$\;
  \While{$a\neq\bot$}{
    $\mathrm{res}\gets\min(\mathrm{res},\,D[v][a]+\mathrm{best}[a])$\;
    $a\gets\mathrm{par}[a]$\;
  }
  \Return $\mathrm{res}$\;
}
\end{algorithm}

\begin{whybox}
\textsc{Mark}$(v)$ walks $v$'s $\Oh(\log n)$ centroid ancestors and tells each
ancestor $a$ ``a marked vertex now sits at distance $D[v][a]$ from you,'' updating
$\mathrm{best}[a]$. \textsc{Query}$(v)$ walks the same ancestors and, for each $a$,
considers routing through $a$: the nearest marked vertex \emph{in $a$'s component}
is $\mathrm{best}[a]$ away from $a$, and $v$ is $D[v][a]$ away from $a$, giving a
candidate path of length $D[v][a]+\mathrm{best}[a]$. Correctness is exactly
Lemma~\ref{ch20:lem-pathkey}: the true nearest marked vertex $s^*$ and $v$ have a
centroid-tree lca $a^*$ lying on the $v$--$s^*$ path, so
$\dist_T(v,s^*)=\dist_T(v,a^*)+\dist_T(a^*,s^*)=D[v][a^*]+\dist_T(a^*,s^*)\ge
D[v][a^*]+\mathrm{best}[a^*]$, and the query examines $a^*$. Every candidate
$D[v][a]+\mathrm{best}[a]$ is the length of an actual walk $v\to a\to s$, hence an
upper bound on $\dist_T(v,s)\ge\min_{s\in S}\dist_T(v,s)$; and the bound is met at
$a=a^*$. So the query returns the exact minimum. Each operation is $\Oh(\log n)$.
\end{whybox}

\begin{trapbox}
The classic mistake is to store, at $\mathrm{best}[c]$, the nearest marked vertex
\emph{in the original subtree of $c$} (an HLD-style rooted-subtree notion) rather
than in $c$'s \emph{centroid component} -- the set of vertices for which $c$ is a
centroid ancestor. Those are different sets, and using the wrong one makes
\textsc{Query} miss valid marked vertices. A second hazard is an inexact distance
table: $D[x][a]$ must be the \emph{true} tree distance, computed by each centroid's
BFS over its live component during the build. If $D$ is right, the candidate
$D[v][a]+\mathrm{best}[a]$ is always a real walk length, so the algorithm cannot
overshoot; the only way to a wrong answer is a wrong $D$.
\end{trapbox}

\begin{labbox}
On our tree, \textsc{Mark}$(1)$. Vertex $1$'s centroid ancestors are
$1\to2\to6$ (parents read off the centroid tree). The true tree distances:
$D[1][1]=0$, $D[1][2]=\dist_T(1,2)=1$ (edge $1$-$2$), and $D[1][6]=\dist_T(1,6)=3$
(path $1\!-\!2\!-\!4\!-\!6$). So after the mark,
$\mathrm{best}[1]=0,\ \mathrm{best}[2]=1,\ \mathrm{best}[6]=3$. Now
\textsc{Query}$(9)$: ancestors $9\to7\to6$ with $D[9][9]=0$, $D[9][7]=1$ (edge
$9$-$7$), $D[9][6]=2$ (path $9\!-\!7\!-\!6$). Candidates:
$D[9][9]+\mathrm{best}[9]=0+\infty$; $D[9][7]+\mathrm{best}[7]=1+\infty$;
$D[9][6]+\mathrm{best}[6]=2+3=5$. The answer $5$ matches the tree path
$9\!-\!7\!-\!6\!-\!4\!-\!2\!-\!1$, which has $5$ edges. The two arms
$9\!-\!7\!-\!6$ and $6\!-\!4\!-\!2\!-\!1$ are vertex-disjoint except at the meeting
centroid $6$, which is precisely why adding $D[9][6]$ and $\mathrm{best}[6]$ gives
the true distance and not an overcount.
\end{labbox}

\section{The centroid tree as a query structure}
\label{ch20:querytree}

Step back and name what we built. The centroid tree is a balanced auxiliary tree
on the same vertex set with three properties worth memorizing.

\begin{enumerate}
  \item \textbf{Logarithmic height} (Theorem~\ref{ch20:thm-depth}): every vertex
        has $\Oh(\log n)$ ancestors, so per-vertex ancestor walks are cheap.
  \item \textbf{Path ownership} (Lemma~\ref{ch20:lem-pathkey}): the lca in the
        centroid tree lies on the original tree path and uniquely owns it.
  \item \textbf{Component nesting}: the components processed down a root-to-vertex
        path of the centroid tree are nested and shrink by at least half, so each
        vertex lives in only $\Oh(\log n)$ components.
\end{enumerate}

These three together yield the universal template: precompute $D[x][a]$
(distance to each centroid ancestor) during the build; keep an aggregate
$\mathrm{agg}[c]$ at each centroid summarizing the marked/active vertices in $c$'s
component; an update walks $\Oh(\log n)$ ancestors pushing information up, a query
walks the same ancestors pulling it down and combining via the path-through-$c$
decomposition. \textsc{Mark}/\textsc{Query} above is the additive instance; with a
Fenwick tree or balanced multiset at each centroid you get range-distance counts,
$k$-th nearest, sum of distances within radius $r$, and so on, all in
$\Oh(\log^2 n)$ per operation.

\begin{applicationbox}
Centroid decomposition is the standard tool behind several recurring problems:
\begin{itemize}
  \item \textbf{``Count paths of length $k$''} and weighted variants (count paths
        with weight $\le W$, using a Fenwick over arm lengths) -- Pattern A.
  \item \textbf{Nearest active facility on a tree under toggles} (telecom or CDN
        placement on a tree-shaped network) -- Pattern B.
  \item \textbf{Sum of pairwise distances} $\sum_{u<v}\dist_T(u,v)$: a single
        pass attributing each path to its centroid.
  \item \textbf{Point updates with path-aggregate queries between arbitrary pairs}
        where the queries are not along a single tree path (so HLD does not apply).
\end{itemize}
The discriminator versus heavy-light decomposition: use HLD when queries follow one
tree path; use centroid decomposition when queries range over many paths or ask
``nearest / all within distance.''
\end{applicationbox}

\begin{connectionbox}
A centroid is the cheapest \emph{vertex separator} of a tree -- one vertex whose
deletion balances the pieces. The planar-separator and treewidth chapters
generalize this: a planar graph has an $\Oh(\sqrt n)$-vertex balanced separator,
and recursing on separators (separator decomposition) is the exact analogue of the
centroid tree, giving nested-dissection orderings for fast linear solves. The
$\Oh(\log n)$-depth recursion tree here is the simplest member of that family. The
centroid tree's logarithmic height also makes it a clean substrate for the
Euler-tour and LCA structures of the next chapters. And when you later see
Laplacian solvers built on recursive vertex elimination, recall that eliminating a
balanced separator first is precisely what keeps fill-in -- and cost -- under
control: centroid decomposition is that idea in its purest, tree-only form.
\end{connectionbox}

\section{Weighted and rooted variants}
\label{ch20:variants}

Two extensions appear constantly and need only small changes.

\textbf{Edge weights.} If edges carry nonnegative lengths and ``distance'' means
summed weight, nothing structural changes: centroids are still defined by
\emph{vertex counts} (so the $\Oh(\log n)$ depth is untouched), and the only edit
is that the BFS from each centroid becomes a weighted DFS accumulating edge lengths
into $D[x][c]$. Pattern A's histogram becomes indexed by weight, often needing a
Fenwick tree to count arms with weight in a range; Pattern B is unchanged.

\textbf{Counting paths with length in a range $[\ell,r]$.} Replace the exact-$k$
convolution with a two-pointer or Fenwick query: for each arm of length $a$, count
opposite arms of length in $[\ell-a,\,r-a]$. Same inclusion--exclusion over
branches.

\begin{formalbox}
\textbf{Why vertex-count balance, not weight balance, is the right invariant.} The
depth bound (Theorem~\ref{ch20:thm-depth}) needs each child component to have at
most half as many \emph{vertices}, because the recursion bottoms out when a
component has one vertex. If you tried to balance by total edge weight, a tree with
one enormously heavy edge could force unbalanced vertex splits and blow the depth
to $\Theta(n)$. Keep the balancing criterion on counts; let weights live only in
the distance values $D[\cdot][\cdot]$.
\end{formalbox}

\section*{Summary}

A centroid splits a tree into pieces of at most half the size; recursing on the
pieces builds the centroid tree, a balanced auxiliary tree of height
$\Oh(\log n)$ computable in $\Oh(n\log n)$. Its defining property is that the
centroid-tree lca of any two vertices lies on -- and uniquely owns -- the tree path
between them. That ownership turns global ``all paths'' questions into per-centroid
local work, and the logarithmic depth keeps the bookkeeping at $\Oh(n\log n)$ for
offline counting (Pattern A) or $\Oh(\log n)$ per operation for online
nearest-marked queries (Pattern B). Reach for it whenever tree queries roam over
many paths rather than tracking a single one.

\section{Exercises}
\label{ch20:exercises}

\begin{exercise}[Implement the build]
Implement Algorithm~\ref{ch20:alg-build} and verify on random trees that the
resulting centroid tree has height $\le\floor{\log_2 n}$. Instrument the code to
assert that every \textsc{SubtreeSizes} call touches only live vertices, and
measure total DFS work to confirm it is $\Oh(n\log n)$, not $\Oh(n^2)$.
\end{exercise}

\begin{exercise}[Two centroids]
Sharpen Theorem~\ref{ch20:thm-exists}: characterize exactly which trees have two
centroids. (Hint: an edge whose removal splits the tree into two halves of size
$n/2$ each; what must $n$ be?) Give the smallest tree with two centroids and the
smallest with a unique centroid, for each parity of $n$.
\end{exercise}

\begin{exercise}[Sum of all pairwise distances]
Use centroid decomposition to compute $\sum_{u<v}\dist_T(u,v)$ in $\Oh(n\log n)$.
Then find an $\Oh(n)$ algorithm using a single tree DP (each edge contributes
$\mathrm{size}\cdot(n-\mathrm{size})$). Explain when each approach is preferable and
why the simpler $\Oh(n)$ method does not generalize to ``distance exactly $k$.''
\end{exercise}

\begin{exercise}[Break the naive size count]
Construct a family of trees on which a centroid-decomposition implementation that
recomputes subtree sizes over the \emph{entire} tree at every level runs in
$\Theta(n^2)$. Confirm empirically, then fix it with the live-component guard and
show the curve flattens to $\Oh(n\log n)$.
\end{exercise}

\begin{exercise}[Paths in a length range]
Extend Algorithm~\ref{ch20:alg-count} to count pairs with $\dist_T(u,v)\in[\ell,r]$
using a Fenwick tree over arm lengths instead of an exact-$k$ convolution. State
and justify the resulting running time, including the extra logarithmic factor.
\end{exercise}

\begin{exercise}[Weighted nearest marked]
Adapt Pattern B (Algorithm~\ref{ch20:alg-nearest}) to nonnegative edge weights.
Prove that the $D[x][a]$ table, built by a weighted traversal from each centroid,
still has only $\Oh(\log n)$ entries per vertex, and that \textsc{Query} remains
exact. Where does the proof use Lemma~\ref{ch20:lem-pathkey}?
\end{exercise}

\begin{exercise}[$k$-th nearest marked vertex]
Augment each centroid with a sorted multiset (or Fenwick over distances) of the
distances from $c$ to marked vertices in its component, so \textsc{Query}$(v)$ can
return the number of marked vertices within distance $r$ of $v$. Achieve
$\Oh(\log^2 n)$ per operation and prove the bound (one logarithm for centroid
ancestors, one for the per-centroid structure). Watch for double counting across
overlapping centroid components and explain how a per-branch correction avoids it.
\end{exercise}

\begin{exercise}[Centroid vs.\ center]
Give a tree where the centroid and the center (the eccentricity minimizer) are
different vertices, and one where they coincide. Conclude with an explicit example
that the center is not a substitute for the centroid in the decomposition: exhibit
a tree whose center, used as the split vertex, produces a child component of size
$>n/2$ and thus a recursion deeper than $\Oh(\log n)$.
\end{exercise}

\begin{exercise}[Why static]
Centroid decomposition as built here is static -- inserting or deleting a tree edge
can change centroids globally. Show that a single edge change can move the top
centroid arbitrarily far, sketch why a fully dynamic centroid decomposition is
hard, and contrast with link-cut trees (later chapters) which \emph{do} support
edge updates. For which of Pattern A / Pattern B does staticness matter, and why?
\end{exercise}

\begin{exercise}[Lower-bound the depth]
Show the height bound $\floor{\log_2 n}$ is essentially tight: exhibit a tree whose
centroid tree has height $\Theta(\log n)$ (e.g.\ a complete binary tree), and argue
no tree forces height $\omega(\log n)$. Then show a \emph{path} graph still yields
centroid-tree height $\Theta(\log n)$ despite being maximally unbalanced as a
rooted tree -- explaining why ``unbalanced as a rooted tree'' and ``deep centroid
decomposition'' are unrelated notions.
\end{exercise}
